\documentclass[10pt,a4paper]{article}
\usepackage{amsmath,amssymb,amsthm,graphicx}
\usepackage[no-math]{fontspec}
\newenvironment{authorclaim}[1]{\par\medskip\noindent\textbf{#1.}\ }{\par\medskip}
\allowdisplaybreaks
\title{Randomized counting-to-decision reduction for exponent-p class-two matrix-group isomorphism}
\author{Joshua A. Grochow, Youming Qiao}
\date{doi:10.4230/LIPIcs.CCC.2021.16; CC BY 4.0}
\begin{document}
\maketitle
\noindent\textit{Editorial scope.} Counted once: randomized group counting-to-decision Theorem B and its full B-prime isometry mechanism. All four gadget requirements, verification, functionality, projective orbit counting, random-subspace/classical-group probability analysis and group multiplicity calculation retained. Preserves all original prime/field/exponent/class/input and randomness conditions; no support results or parameter regimes counted separately.\par\medskip
\noindent\textit{Allowed standard prerequisites (editorial).}\begin{itemize}
\item Finite-field subspace counting, matrix-algebra/module algorithms and random generation of classical groups. Complete source-local Section5 proof retained, including Claims 19 and20 and Proposition15 probability analysis. Exact imported Theorem21 on generation of SL(n,q) and constructive Baer group transfer retained with the complete74-entry bibliography. Imported classical group and module algorithms are not recursively reproved.
\end{itemize}\medskip
% Faithful formal-publication context for search/count reductions: Sections1.2 and2.
\section*{Author notation and isometry models}
\paragraph{Vector spaces.} Let $\mathbb F$ be a field. In this paper we only consider finite-dimensional vector spaces over $\mathbb F$. We use $\mathbb F^n$ to denote the vector space of length-$n$ column vectors. The $i$th standard basis vector of $\mathbb F^n$ is denoted $\vec e_i$. Depending on the context, 0 may denote the zero vector space, a zero vector, or an all-zero matrix. For $S$ a set of vectors, we use $\langle S\rangle$ to denote the subspace spanned by elements in $S$.
\paragraph{Some groups.} The general linear group of degree $n$ over a field $\mathbb F$ is denoted by $\mathrm{GL}(n,\mathbb F)$. The symmetric group of degree $n$ is denoted by $S_n$. The natural embedding of $S_n$ into $\mathrm{GL}(n,\mathbb F)$ is to represent permutations by permutation matrices. The subgroup of $\mathrm{GL}(n,\mathbb F)$ consisting of diagonal matrices is called the diagonal subgroup, denoted by $\operatorname{diag}(n,\mathbb F)$. A monomial matrix is a product of a diagonal and a permutation matrix; equivalently, each row and each column has exactly one non-zero entry. The collection of monomial matrices forms a subgroup of $\mathrm{GL}(n,\mathbb F)$, which we call the monomial subgroup and denote by $\operatorname{Mon}(n,\mathbb F)$. It is the semi-direct product $\operatorname{diag}(n,\mathbb F)\rtimes S_n\cong(\mathbb F^*)^n\rtimes S_n$.
\paragraph{Matrices.} Let $\mathrm M(\ell\times n,\mathbb F)$ be the linear space of $\ell\times n$ matrices over $\mathbb F$, and $\mathrm M(n,\mathbb F):=\mathrm M(n\times n,\mathbb F)$. Given $A\in\mathrm M(\ell\times n,\mathbb F)$, $A^t$ denotes the transpose of $A$.
A matrix $A\in\mathrm M(n,\mathbb F)$ is alternating, if for any $u\in\mathbb F^n$, $u^tAu=0$. That is, $A$ represents an alternating bilinear form. Note that in characteristic $\ne2$, alternating is the same as skew-symmetric, but in characteristic 2 they differ (in characteristic 2, skew-symmetric=symmetric). The linear space of $n\times n$ alternating matrices over $\mathbb F$ is denoted by $\Lambda(n,\mathbb F)$.
The $n\times n$ identity matrix is denoted by $I_n$, and when $n$ is clear from the context, we may just write $I$. The elementary matrix $E_{i,j}$ is the matrix with the $(i,j)$th entry being 1, and other entries being 0. The $(i,j)$-th elementary alternating matrix is the matrix $E_{i,j}-E_{j,i}$.
\paragraph{Matrix tuples.} We use $\mathrm M(\ell\times n,\mathbb F)^m$ to denote the linear space of $m$-tuples of $\ell\times n$ matrices. Boldface letters like $\mathbf A$ and $\mathbf B$ denote matrix tuples. Let $\mathbf A=(A_1,\ldots,A_m)$, $\mathbf B=(B_1,\ldots,B_m)\in\mathrm M(\ell\times n,\mathbb F)^m$. Given $P\in\mathrm M(\ell,\mathbb F)$ and $Q\in\mathrm M(n,\mathbb F)$, $P\mathbf A Q:=(PA_1Q,\ldots,PA_mQ)\in\mathrm M(\ell,\mathbb F)$. Given $R=(r_{i,j})_{i,j\in[m]}\in\mathrm M(m,\mathbb F)$, $\mathbf A^R:=(A'_1,\ldots,A'_m)\in\mathrm M(m,\mathbb F)$ where $A'_i=\sum_{j\in[m]}r_{j,i}A_j$.
\begin{authorclaim}{Remark5}
In particular, note that the coefficients in the formula of defining $A'_i$ correspond to the entries in the $i$th column of $R$. While this choice is immaterial (we could have chosen the opposite convention), all of our later calculations are consistent with this convention.
\end{authorclaim}
Given $\mathbf A,\mathbf B\in\mathrm M(\ell\times n,\mathbb F)^m$, we say that $\mathbf A$ and $\mathbf B$ are isometric, if there exists $P\in\mathrm{GL}(n,\mathbb F)$, such that $P^t\mathbf A P=\mathbf B$. Finally, $\mathbf A$ and $\mathbf B$ are pseudo-isometric if there exist $P\in\mathrm{GL}(n,\mathbb F)$ and $R\in\mathrm{GL}(m,\mathbb F)$, such that $P^t\mathbf A P=\mathbf B^R$.
\paragraph{Matrix spaces.} Linear subspaces of $\mathrm M(\ell\times n,\mathbb F)$ are called matrix spaces. Calligraphic letters like $\mathcal A$ and $\mathcal B$ denote matrix spaces. By a slight abuse of notation, for $\mathbf A\in\mathrm M(\ell\times n,\mathbb F)^m$, we use $\langle\mathbf A\rangle$ to denote the subspace spanned by those matrices in $\mathbf A$. For $\mathbf A,\mathbf B\in\mathrm M(n,\mathbb F)^m$, we say that the spaces $\langle\mathbf A\rangle,\langle\mathbf B\rangle$ are isometric iff the tuples $\mathbf A,\mathbf B$ are pseudo-isometric.
\begin{authorclaim}{Definition4 (The Alternating Matrix Space Isometry problem)}
Input: $A_1,\ldots,A_m$ and $B_1,\ldots,B_m$, $n\times n$ alternating matrices over a field $\mathbb F$.
Decide: Is there a $L\in\mathrm{GL}(n,\mathbb F)$, such that the linear span of $\{A_i:i\in[m]\}$ is equal to the linear span of $\{L^tB_iL:i\in[m]\}$?
\end{authorclaim}
\paragraph{3-way arrays.} Let $\mathrm T(\ell\times n\times m,\mathbb F)$ be the linear space of $\ell\times n\times m$ 3-way arrays over $\mathbb F$. We use the fixed-width teletypefont for 3-way arrays, like $\mathtt A,\mathtt B$, etc..
Given $\mathtt A\in\mathrm T(\ell\times n\times m,\mathbb F)$, we can think of $\mathtt A$ as a 3-dimensional table, where the $(i,j,k)$th entry is denoted as $\mathtt A(i,j,k)\in\mathbb F$. We can slice $\mathtt A$ along one direction and obtain several matrices, which are then called slices. For example, slicing along the first coordinate, we obtain the horizontal slices, namely $\ell$ matrices $A_1,\ldots,A_\ell\in\mathrm M(n\times m,\mathbb F)$, where $A_i(j,k)=\mathtt A(i,j,k)$. Similarly, we also obtain the lateral slices by slicing along the second coordinate, and the frontal slices by slicing along the third coordinate.
We will often represent a 3-way array as a matrix whose entries are vectors. That is, given $\mathtt A\in\mathrm T(\ell\times n\times m,\mathbb F)$, we can write
\[\mathtt A=\begin{bmatrix}w_{1,1}&w_{1,2}&\cdots&w_{1,n}\\w_{2,1}&w_{2,2}&\cdots&w_{2,n}\\\vdots&\vdots&\ddots&\vdots\\w_{\ell,1}&w_{\ell,2}&\cdots&w_{\ell,n}\end{bmatrix},\]
where $w_{i,j}\in\mathbb F^m$, so that $w_{i,j}(k)=\mathtt A(i,j,k)$. Note that, while $w_{i,j}\in\mathbb F^m$ are column vectors, in the above representation of $\mathtt A$, we should think of them as along the direction ``orthogonal to the paper.'' Following [45], we call $w_{i,j}$ the tube fibers of $\mathtt A$. Similarly, we can have the row fibers $v_{i,k}\in\mathbb F^n$ such that $v_{i,k}(j)=\mathtt A(i,j,k)$, and the column fibers $u_{j,k}\in\mathbb F^\ell$ such that $u_{j,k}(i)=\mathtt A(i,j,k)$.
Given $P\in\mathrm M(\ell,\mathbb F)$ and $Q\in\mathrm M(n,\mathbb F)$, let $P\mathtt A Q$ be the $\ell\times n\times m$ 3-way array whose $k$th frontal slice is $PA_kQ$. For $R=(r_{i,j})\in\mathrm{GL}(m,\mathbb F)$, let $\mathtt A^R$ be the $\ell\times n\times m$ 3-way array whose $k$th frontal slice is $\sum_{k'\in[m]}r_{k',k}A_{k'}$. Note that these notations are consistent with the notations for matrix tuples above, when we consider the matrix tuple $\mathbf A=(A_1,\ldots,A_k)$ of frontal slices of $\mathtt A$.

\begin{authorclaim}{TheoremB}
For $p$ an odd prime, $p\ge n^{\Omega(1)}$, there is a randomized counting-to-decision reduction for $\operatorname{GpIso2Exp}(p)$ for groups of order $p^n$, running in time $p^{O(n)}=\operatorname{poly}(|G|)$.
\end{authorclaim}

\begin{authorclaim}{TheoremB'}
For $q$ a prime power with $q=n^{\Omega(1)}$, there is a randomized counting-to-decision reduction for Alternating Matrix Space Isometry which, given $n\times n$ alternating matrix spaces $\mathcal A,\mathcal B$ over $\mathbb F_q$ of dimension $m$, computes the number of isometries from $\mathcal A$ to $\mathcal B$ in time $q^{O(n)}$. The reduction queries the decision oracle with inputs of dimension at most $O(n^2)$.
\end{authorclaim}

\section*{5 Counting-to-decision reduction by restricting to diagonal groups}
In this section, we devise a gadget to achieve the restriction to the group of diagonal matrices, and use it to do the counting to decision reduction for Alternating Matrix Space Isometry.
\subsection*{5.1 Preliminaries}
Some preparations are in order.
\begin{authorclaim}{Observation12}
Let $n\ge23$. Then any permutation $\sigma\in S_n$ either fixes a set of 6 points $P\subseteq[n]$, or moves a set of 6 points $P\subseteq[n]$ to another set of 6 points $Q\subseteq[n]$ such that these two sets are disjoint.
\end{authorclaim}
\begin{proof}
Suppose $\sigma$ fixes at most 5 points. Then there are at least 18 points that are not fixed by $\sigma$. Suppose $\sigma$ has $t$ non-trivial cycles of length $l_1,\ldots,l_t$, such that $\sum_i l_i\ge18$. For a cycle $(p_1,\ldots,p_s)$, we can choose $p_1,p_3,\ldots,p_{2\lfloor s/2\rfloor-1}$ and put them in $P$, and $p_2,p_4,\ldots,p_{2\lfloor s/2\rfloor}$ in $Q$. Do this for every cycle, we obtain the desired $P$ and $Q$. The worst case is when every cycle is of length 3. Since there are at least 18 points not fixed by $\sigma$, $P$ is of size $\ge6$.
\end{proof}
We shall make repeated uses of the following facts.
\begin{authorclaim}{Fact13}
\begin{enumerate}
\item Given $a_i\in\mathbb R$, $0\le a_i\le1$, $i\in[m]$, $\prod_{i\in[m]}(1-a_i)\ge1-\sum_{i\in[m]}a_i$.
\item Let $m,N\in\mathbb N$ and $1\le m\le N$. A random matrix $A\in\mathrm M(N\times m,q)$ is of rank $m$ with probability $\ge1-2/q^{N-m+1}$.
\item For $n\in\mathbb N$, $0\le d\le n$, the number of dimension-$d$ subspaces of $\mathbb F_q^n$ is equal to the Gaussian binomial coefficient
\[\binom nd_q:=\frac{(q^n-1)\cdot(q^n-q)\cdots(q^n-q^{d-1})}{(q^d-1)\cdot(q^d-q)\cdots(q^d-q^{d-1})}.\]
\item The Gaussian binomial coefficient satisfies $q^{(n-d)d}\le\binom nd_q\le q^{(n-d)d+d}$.
\item For $d\in\mathbb N$, the number of complement subspaces of a fixed dimension-$d$ subspace of $\mathbb F_q^n$ is $q^{d(n-d)}$.
\end{enumerate}
\end{authorclaim}
\begin{proof}
For(2), $\Pr[\operatorname{rk}(A)=m]=(1-1/q^N)\cdot(1-q/q^N)\cdots(1-q^{m-1}/q^N)$. By(1), we have
\[\Pr[\operatorname{rk}(A)=m]\ge1-\sum_{i=N-m+1}^N\frac1{q^i}=1-\frac1{q^{N-m+1}}-\sum_{i=N-m+2}^N\frac1{q^i}\ge1-\frac2{q^{N-m+1}}.\]
\end{proof}
\subsection*{5.2 Describing the gadget}
Let $\mathcal A\le\Lambda(n,q)$ be an alternating matrix space, and let $\mathbf A=(A_1,\ldots,A_m)\in\Lambda(n,q)^m$ be an ordered linear basis of $\mathcal A$. Let $\mathtt A\in\mathrm T(n\times n\times m,\mathbb F_q)$ be the 3-way array constructed from $\mathbf A$, i.e. the $i$th frontal slice of $\mathtt A$ is $A_i$.
We shall assume $n=\Omega(1)$, and $q=n^{\Omega(1)}$ throughout the remainder of this section.
\paragraph{The form of the gadget.} To describe the gadget, it is easier to view $\mathtt A$ from the lateral viewpoint. That is, for $i\in[n]$, let $C_i=[A_1e_i,\ldots,A_me_i]\in\mathrm M(n\times m,q)$. Let $\mathbf C=(C_1,\ldots,C_n)\in\mathrm M(n\times m,q)^n$. Then construct $\mathbf C'=(C'_1,\ldots,C'_n)$, $C'_i=\begin{bmatrix}C_i&0\\0&G_i\end{bmatrix}$, where $G_i$ is of size $6n\times4n^2$. For $i\in[n]$, $G_i=[0\ \cdots\ 0\ H_i\ 0\ \cdots\ 0]$, where $H_i$ is of size $6n\times4n$ in the $i$th block, and 0 denotes an all-zero matrix of size $6n\times4n$. The $H_i$ will be described below.
Note that from the frontal viewpoint of looking at $\mathtt A$, $G_i$'s are inserted, vertically, below and behind $\mathtt A$. So to preserve the alternating structure, $-G_i$'s also need to be inserted, horizontally, on the right and behind $\mathtt A$. We therefore get $\widetilde{\mathtt A}$, which is of size $7n\times7n\times(m+4n^2)$.
\paragraph{Conditions imposed on the $H_i$'s.} Of course, the key to the construction above lies in the properties of the $H_i$'s. Let $V_i\le\mathbb F_q^{6n}$ be the subspace spanned by the columns of $H_i$. We shall impose the following conditions on $H_i$.
\begin{enumerate}
\item For any $i\in[n]$, $\operatorname{rk}(H_i)=\dim(V_i)=4n$.
\item For any $i,j\in[n]$, $i\ne j$, $\operatorname{rk}([H_i\ H_j])=\dim(V_i\cup V_j)=6n$.
\item For any $(i_1,i_2,i_3,i_4,i_5,i_6)\in[n]^6$ and $(j_1,j_2,j_3,j_4,j_5,j_6)\in[n]^6$, such that $|\{i_1,\ldots,i_6\}\cup\{j_1,\ldots,j_6\}|=12$, i.e. $i_k$ and $j_\ell$ all different, the coset $\mathcal C=\{T\in\mathrm{GL}(6n,q):\forall k\in[6],T(V_{i_k})=V_{j_k}\}$ is empty. Note that for any $i\in[n]$, $T(V_i)$ is spanned by the columns of $TH_i$.
\item For any $(i_1,i_2,i_3,i_4,i_5,i_6)\in[n]^6$, $i_k$ all different, the group $\mathcal S=\{T\in\mathrm{GL}(6n,q):\forall k\in[6],T(V_{i_k})=V_{i_k}\}$ consists of only of scalar matrices.
\end{enumerate}
\begin{authorclaim}{Remark14}
Given $H_1,\ldots,H_n\in\mathrm M(6n\times4n,q)$, whether they satisfy the four conditions can be verified in polynomial time.
\end{authorclaim}
Conditions(1) and(2) are easily verified in deterministic polynomial time.
For condition(3), it can be formulated as a linear algebraic problem as follows. Let $X$ be a $6n\times6n$ variable matrix. Let $Y_k$, $k\in[6]$, be $4n\times4n$ variable matrices. Set up the equations $XH_{i_k}=H_{j_k}Y_k$, and solve the linear equations to get a subspace of $\mathbb F_q^{(6n)^2+6\cdot(4n)^2}$. The question is then whether this subspace contains $(T,R_1,\ldots,R_6)$ where $T\in\mathrm{GL}(6n,q)$ and $R_i\in\mathrm{GL}(4n,q)$. This is an instance of the symbolic determinant identity testing (SDIT) problem, so it admits a randomized efficient algorithm when $q=n^{\Omega(1)}$.
In fact, this instance of SDIT problem can be solved in deterministic polynomial time. For this let us also check out condition(4). Here, let $X$ and $Y_i$ be from above, and set up the equations $XH_{i_k}=H_{i_k}Y_k$. Solve the linear equations to get a subspace of $\mathbb F_q^{(6n)^2+6\cdot(4n)^2}$. This subspace turns out to be an algebra under the natural multiplications. Indeed, if $AH_{i_k}=H_{i_k}B_k$ and $A'H_{i_k}=H_{i_k}B'_k$, then $AA'H_{i_k}=H_{i_k}B_kB'_k$. To compute the unit group in a matrix algebra can be solved by a polynomial-time Las Vegas algorithm by [16]. Given the unit group, whether it consists of only scalar matrices can be verified easily in deterministic polynomial time.
Then the linear space in condition(3) is a module over the algebra defined in the last paragraph. Because of this structure, the SDIT problem for such instances can be solved in deterministic polynomial time [14,19,37].
\subsection*{5.3 Construction and properties of the gadget}
The following three propositions reveal the construction and functions of the gadget described above. First about the construction. Instead of constructing the above $H_i$'s explicitly in a deterministic way, we shall show that random choices suffice.
\begin{authorclaim}{Proposition15}
Let $H_i\in\mathrm M(6n\times4n,q)$, $i\in[n]$, be random matrices. Then $H_i$'s satisfy the four conditions in Section5.2 with probability $\ge1-n^{O(1)}/q^{\Omega(1)}$.
\end{authorclaim}
Second about the functionality. The following proposition formally explains this.
\begin{authorclaim}{Proposition16}
Suppose $\mathtt A$ and $\mathtt B$ are two 3-tensors constructed from ordered bases of $m$-dimensional alternating matrix spaces $\mathcal A,\mathcal B\le\Lambda(n,q)$. Let $\widetilde{\mathtt A}$ and $\widetilde{\mathtt B}$ be constructed as above, and let $\widetilde{\mathcal A}$ and $\widetilde{\mathcal B}$ be the alternating matrix spaces spanned by the frontal slices of $\widetilde{\mathtt A}$ and $\widetilde{\mathtt B}$, respectively. Then $\mathcal A$ and $\mathcal B$ are isometric via a diagonal matrix if and only if $\widetilde{\mathcal A}$ and $\widetilde{\mathcal B}$ are isometric.
\end{authorclaim}
Finally we shall use this gadget to achieve a counting-to-decision reduction for Alternating Matrix Space Isometry. Formally, we have the following.
\begin{authorclaim}{Proposition17}
Suppose we are given $\mathcal A,\mathcal B\le\Lambda(n,q)$ and a decision oracle for Alternating Matrix Space Isometry. Then there exists a Las Vegas randomized algorithm that computes the number of isometries from $\mathcal A$ to $\mathcal B$ in time $q^{O(n)}$.
\end{authorclaim}
The next three subsections are devoted to the proofs of Propositions15 (Section5.3.3),16 (Section5.3.1), and17 (Section5.3.2). Note that, because the proof of Proposition15 is more complicated compared to the other two, we postpone it to the last.
\begin{authorclaim}{Remark18}
In fact, we expect that this construction works even for small finite fields. The bottleneck lies in Proposition15. If the probability $n^{O(1)}/q^{\Omega(1)}$ could be improved to $n^{O(1)}/q^{\Omega(n)}$, then we would be done. We believe it possible to utilize the structure of invariant subspaces under matrix actions over $\mathbb F_q$ to achieve this. However, we expect that the calculations will be tedious and heavy, so we hope to leave this to a future work.
\end{authorclaim}
\subsection*{5.3.1 Restricting to the diagonal group}
Briefly speaking, conditions 1 and 2 ensure that we first restrict to monomial matrices. Conditions 3 and 4 prevent non-trivial permutations due to the following. As we assume $n=\Omega(1)$, by Observation12, $\sigma\in S_n$ either fixes 6 elements in $[n]$, or moves a set of 6 elements to another, disjoint, set of 6 elements. Condition 3 ensures that the second case could not happen. Condition 4 ensures that in the first case, the only possible invertible matrices that ``preserves'' the matrices $G_i$ for $i\in P$ when multiplying from the left are scalar matrices.
We now prove Proposition16.
\begin{proof}[Proof of Proposition16]
Recall that we construct such $\widetilde{\mathtt A}$ and $\widetilde{\mathtt B}$ from $\mathtt A$ and $\mathtt B$, respectively, using the method in Section5.2. Let $\widetilde{\mathcal A}$ and $\widetilde{\mathcal B}$ be alternating matrix spaces in $\Lambda(7n,q)$, spanned by the frontal slices of $\widetilde{\mathtt A}$ and $\widetilde{\mathtt B}$, respectively.
We want to show that $\widetilde{\mathcal A}$ and $\widetilde{\mathcal B}$ are isometric if and only if $\mathcal A$ and $\mathcal B$ are isometric via diagonal matrices. The if direction is straightforward. Suppose there exist $P=\operatorname{diag}(\alpha_1,\ldots,\alpha_n)\in\operatorname{diag}(n,q)$ and $Q\in\mathrm{GL}(m,q)$ such that $P^t\mathbf A P=\mathbf B^Q$. Let $\widetilde P=\begin{bmatrix}P&0\\0&I_{6n}\end{bmatrix}\in\mathrm{GL}(7n,q)$. Let $\widetilde Q=\begin{bmatrix}Q&0\\0&Q'\end{bmatrix}\in\mathrm{GL}(m+4n^2)$, where $Q'=\operatorname{diag}(\alpha_1I_{4n},\ldots,\alpha_nI_{4n})$. Then it is easy to verify that $\widetilde P^t\widetilde{\mathtt A}\widetilde P=\widetilde{\mathtt B}^{\widetilde Q}$.
Now we turn to the only if direction. If $\widetilde{\mathcal A}$ and $\widetilde{\mathcal B}$ are isometric, then there exists $\widetilde P\in\mathrm{GL}(7n,q)$ and $\widetilde Q\in\mathrm{GL}(m+4n^2,q)$, such that $\widetilde P^t\widetilde{\mathtt A}\widetilde P=\widetilde{\mathtt B}^{\widetilde Q}$. Let $\widetilde P=\begin{bmatrix}P_{1,1}&P_{1,2}\\P_{2,1}&P_{2,2}\end{bmatrix}$, where $P_{1,1}$ is of size $n\times n$. It can be checked easily, from the lateral viewpoint, that $P_{1,2}=0$. As if not, then some $H_i$ would appear in one of the last $6n$ lateral slices in $\widetilde{\mathtt A}^{\widetilde P}$. This would set this slice to be of rank $\ge4n$ by condition(1), which contradicts that the corresponding lateral slice of $\widetilde{\mathtt B}^{\widetilde Q}$ is of rank $\le n$. It follows that $P_{1,1}\in\mathrm{GL}(n,q)$ and $P_{2,2}\in\mathrm{GL}(6n,q)$.
We first claim that $P_{1,1}$ has to be a monomial matrix. If not, then one of the first $n$ lateral slice of $\widetilde{\mathtt A}^{\widetilde P}$ has two distinct $H_i$ and $H_j$. By condition(2), this slice is of rank $\ge6n$, which contradicts that the corresponding lateral slice of $\widetilde{\mathtt B}^{\widetilde Q}$ is of rank $\le5n$.
We further claim that $P_{1,1}$ has to be a diagonal matrix. If not, then suppose the non-trivial permutation underlying $P_{1,1}$ is $\sigma\in S_n$. Since we assumed $n=\Omega(1)$, by Observation12, one of the following two cases has to happen.
\begin{itemize}
\item $\exists\{i_1,\ldots,i_6\}\subseteq[n]$, $\{j_1,\ldots,j_6\}\subseteq[n]$, $|\{i_1,\ldots,i_6\}\cup\{j_1,\ldots,j_6\}|=12$, such that $\sigma(i_k)=j_k$ for $k\in[6]$. We then claim the following.
\begin{authorclaim}{Claim19}
For $\widetilde P^t\widetilde{\mathtt A}\widetilde P=\widetilde{\mathtt B}^{\widetilde Q}$ to hold, a necessary condition is that $\forall k\in[6]$, $P_{2,2}H_{j_k}$ and $H_{i_k}$ have the same linear span.
\end{authorclaim}
\begin{proof}
To see this, note that the $i_k$th lateral slice of $\widetilde P^t\widetilde{\mathtt A}\widetilde P$ is the $j_k$th lateral slice of $\widetilde P^t\widetilde{\mathtt A}$ (up to a scalar multiple). It is equal to the $i_k$th lateral slice of $\widetilde{\mathtt B}^{\widetilde Q}$. Then $\widetilde P^t$ acts on the left on the $j_k$th lateral slice of $\widetilde{\mathtt A}$. Noting that $P^t=\begin{bmatrix}P_{1,1}^t&P_{2,1}^t\\0&P_{2,2}^t\end{bmatrix}$ and the $j_k$th lateral slice of $\widetilde{\mathtt A}$ is $C'_{j_k}=\begin{bmatrix}C_{j_k}&0\\0&G_{j_k}\end{bmatrix}$, we see that $P^tC'_{j_k}=\begin{bmatrix}*&*\\0&P_{2,2}^tG_{j_k}\end{bmatrix}$. (Here, $C_i$ and $G_i$ are defined in Section5.2.) On the other hand, we see that the $i_k$th lateral slice of $\widetilde{\mathtt B}^{\widetilde Q}$ is the $i_k$th lateral slice multiplied from the right by $\widetilde Q$. Our claim follows then by comparing the last $6n$ rows.
\end{proof}
But the condition(3) excludes the existence of such $P_{2,2}$, so this cannot happen.
\item $\exists\{i_1,\ldots,i_6\}\subseteq[n]$, $i_k$ all different, such that $\sigma(i_k)=i_k$. In this case, for $\widetilde P^t\widetilde{\mathtt A}\widetilde P=\widetilde{\mathtt B}^{\widetilde Q}$ to hold, by the same argument as in the proof of Claim19, a necessary condition is that $P_{2,2}H_{i_k}$ and $H_{i_k}$ have the same linear span. Then the condition(4) ensures that $P_{2,2}=\lambda I_{6n}$ for some $\lambda\ne0\in\mathbb F$ in this setting. Then because $\sigma$ is non-trivial, $\sigma$ moves some $i\in[n]$ to $j\in[n]$, $i\ne j$. By comparing the $j$th lateral slice of $\widetilde P^t\widetilde{\mathtt A}$ and the $i$th lateral slice of $\widetilde{\mathtt B}^{\widetilde Q}$, $P_{2,2}H_i=\lambda H_i$ and $H_j$ have the same linear span, which is not possible because the condition(2) ensures that $H_i$ and $H_j$ span different subspaces.
\end{itemize}
We then have shown that $P_{1,1}$ must be a diagonal matrix. By comparing the top-left-front sub-tensors of size $n\times n\times m$ of $\widetilde P^t\widetilde{\mathtt A}\widetilde P$ and $\widetilde{\mathtt B}^{\widetilde Q}$, we arrive at the desired conclusion that $\mathcal A$ and $\mathcal B$ are isometric via the diagonal matrix $P_{1,1}$.
\end{proof}

\subsection*{5.3.2 Using the gadget for counting-to-decision reduction}
The strategy follows closely the counting to decision reduction for graph isomorphism. We then apply the above strategy to get a counting to decision reduction for alternating matrix space isometry to prove Proposition17.
\begin{proof}[Proof of Proposition17]
Our goal is to compute the number of isomorphisms from $\mathcal A$ to $\mathcal B$, where $\mathcal A,\mathcal B\le\Lambda(n,q)$ are of dimension $m$. First, we use the decision oracle first to decide whether $\mathcal A$ and $\mathcal B$ are isometric. If not, the number of isometries is 0. If so, we need to caculate the order of the autometry group of $\mathcal A$, $\operatorname{Aut}(\mathcal A)$. To do that, we first randomly sample $n$ $6n\times4n$ matrices $H_1,\ldots,H_n$ over $\mathbb F_q$, and verify whether they satisfy the four conditions in Section5.2 using Remark14. Note that this is where the algorithm needs to be a Las Vegas algorithm.
Let $A=\operatorname{Aut}(\mathcal A)$. Recall that $e_i$ denotes the $i$th standard basis vector in $\mathbb F_q^n$. For $i\in[n]$, let $A_i=\{T\in A:\forall1\le j\le i,T(e_i)=\lambda_ie_i,\lambda_i\ne0\in\mathbb F_q\}$. Note that $A_n=A\cap\operatorname{diag}(n,q)$. We can calculate the order of $A_n$ in time $q^{O(n)}$ by brute-force, i.e., enumerating all invertible diagonal matrices. Set $A_0=A$. We then have the tower of subgroups $A_0\ge A_1\ge\cdots\ge A_n$.
To compute the order of $A_0$, it is enough to compute $[A_i:A_{i+1}]$. Note that for $T,T'\in A_i$, $TA_{i+1}=T'A_{i+1}$ as left cosets in $A_i$ if and only if $T(e_{i+1})=\lambda T'(e_{i+1})$ for some $\lambda\ne0\in\mathbb F_q$. So $[A_i:A_{i+1}]$ is equal to the size of the orbit of $e_{i+1}$ under $A_i$ in the projective space. Let $v\in\mathbb F_q^n$. To test whether $v$ is in the orbit of $e_{i+1}$ under $A_i$ in the projective space, we tranform $\mathcal A$ by $P^t\cdot P$, where $P\in\mathrm{GL}(n,q)$ sends $e_{i+1}$ to $v$ and $e_j$ to $e_j$ for $j\ne i+1$, to get $\mathcal A'$. We then add the diagonal restriction gadget to the first $i+1$ lateral slices and the first $i+1$ horizontal slices of $\mathcal A$ and $\mathcal A'$, to obtain $\widetilde{\mathcal A}$ and $\widetilde{\mathcal A}'$ respectively. Then feed $\mathcal A$ and $\mathcal A'$ to the decision oracle. By the functionality of the diagonal restriction gadget, $v$ is in the orbit of $e_{i+1}$ in the projective space if and only if $\widetilde{\mathcal A}$ and $\widetilde{\mathcal A}'$ are isometric. Enumerating $v\in\mathbb F_q^n$ up to scalar multiples gives us the size of the orbit of $e_{i+1}$ under $A_i$ in the projective space. This finishes the description of the algorithm.
A small caveat in the above is that our gadget requires $n=\Omega(1)$, so we cannot start from $A_0$ at the beginning. This issue can be revolved by noting that the order of $A_c$, for any constant $c$, can be computed in time $q^{O(n)}$, by enumerating all possible images of $e_1,\ldots,e_c$ in time $q^{O(n)}$, adding the diagonal restriction gadget, and utilizing the decision oracle.
\end{proof}
\subsection*{5.3.3 Random $H_i$'s satisfy the requirements when $q=n^{\Omega(1)}$}
In the following we will encounter random matrices over $\mathbb F_q$ as well as random subspaces in $\mathbb F_q^n$. There is a subtle point which we want to clarify now. Let $m\le n$. Note that there are $\binom nm_q$ subspaces of $\mathbb F_q^n$, and there are $N_1=(q^n-1)\cdots(q^n-q^{m-1})$ rank-$m$ matrices of size $n\times m$. It can be seen easily that each $m$-dimensional subspace $V$ of $\mathbb F_q^n$ has $N_2=(q^m-1)\cdots(q^m-q^{m-1})$ many representations as rank-$m$ matrices of size $n\times m$, i.e. the columns of the matrix span $V$. It follows that we can work with random rank-$m$ matrices of size $n\times m$ as if we are working with random $m$-dimensional subspaces of $\mathbb F_q^n$. Such correspondences will be used implicitly for other structures, including direct sum decompositions.
Now let us get back to our question. We shall show that a random choice of $H_i$, $i\in[n]$, would satisfy the four conditions we imposed on $H_i$'s. We will prove that for conditions $k=1,2,3$,
\[\Pr[\text{random }H_i\text{ not satisfy condition }k]\le\frac{n^{O(1)}}{q^{\Omega(n)}}.\]
Once these hold, by a union bound, we have
\[\Pr[\exists i\in[3],\text{ random }H_i\text{ not satisfy condition }i]\le\frac{n^{O(1)}}{q^{\Omega(n)}}.\]
For condition(4), we will prove that
\[\Pr[\text{random }H_i\text{ not satisfy condition 4}\mid H_i\text{ satisfy conditions 1,2,3}]\le\frac{n^{O(1)}}{q^{\Omega(1)}}.\]
This then would allow us to conclude that when $q=n^{\Omega(1)}$, random $H_i$'s satisfy all the four conditions. We examine the first three conditions one by one.
\begin{enumerate}
\item For condition(1), by Fact13(2), we have $\Pr[\exists i\in[n],\operatorname{rk}(H_i)<4n]\le n\cdot\Pr[\operatorname{rk}(H_i)<4n]\le2n/q^{2n+1}$.
\item For condition(2), noting that the block matrix $(H_i\ H_j)$ is a random $6n\times8n$ matrix over $\mathbb F_q$, by Fact13(2), we have
\[\Pr[\exists i\ne j\in[n],\operatorname{rk}((H_i\ H_j))<6n]\le\binom n2\cdot\frac2{q^{8n-6n+1}}\le\frac{n^2}{q^{2n+1}}.\]
\item For condition(3), let $I=(H_{i_1}\ \cdots\ H_{i_6})$, and $J=(H_{j_1}\ \cdots\ H_{j_6})$. We see that $\mathcal C$ is non-empty if and only if there exists $L\in\mathrm{GL}(6n,q)$ and $R_k\in\mathrm{GL}(4n,q)$, $k\in[6]$, such that $LH_{i_k}R_k=H_{j_k}$. Note that the orbit of $I$ under this group action is of size at most $q^{(6n)^2+6\cdot(4n)^2}=q^{132n^2}$. Since $i_k$ and $j_\ell$ are all different, the probability of $J$ belonging to this orbit is $\le q^{132n^2}/q^{144n^2}=1/q^{12n^2}$. We then have
\[\Pr[\exists i_k,j_k\in[n],k\in[6],i_k,j_k\text{ all different},\mathcal C=\emptyset]\le\binom n{12}\frac2{q^{12n^2}}\le\frac{n^{12}}{q^{12n^2}}.\]
\end{enumerate}
We now focus on condition(4). For condition(4), we first assume that the conditions(1) and(2) as above hold. Then $V_i$'s are random $4n$-dimensional subspaces of $\mathbb F_q^{6n}$. Note that
\[\Pr[\exists i_k\in[n],k\in[6],i_k\text{ all different},\mathcal S\text{ non-scalar}]\le n^6\cdot\Pr[\mathcal S\text{ non-scalar for }V_1,\ldots,V_6].\]
So we turn to study $\Pr[\mathcal S\text{ non-scalar for }V_1,\ldots,V_6]$, and will show that it is $\le1/q^{\Omega(1)}$.
Let $U_1=V_1\cap V_2$, $U_2=V_2\cap V_3$, and $U_3=V_1\cap V_3$. Let $W_1=V_4\cap V_5$, $W_2=V_5\cap V_6$, and $W_3=V_4\cap V_6$. Since conditions(1) and(2) hold, we have $\dim(U_i)=\dim(W_i)=2n$. We claim that with probability $\ge1-2/q$, $\mathbb F_q^{6n}=U_1\oplus U_2\oplus U_3$, i.e., $U_1\cup U_2\cup U_3$ span $\mathbb F_q^{6n}$. This can be seen as follows. Since we assumed conditions(1) and(2), this happens if and only if $V_1\cap V_2$ and $V_3$ together span $\mathbb F_q^{6n}$. Therefore we calculate, using Fact13(1),(3), and(5), that
\begin{align*}
&\Pr[V_3\text{ is a complement subspace of }V_1\cap V_2]\\
&=q^{2n\cdot4n}/\binom{6n}{4n}_q
=\frac{(q^{6n}-q^{2n})(q^{6n}-q^{2n+1})\cdots(q^{6n}-q^{6n-1})}{(q^{6n}-1)(q^{6n}-q)\cdots(q^{6n}-q^{4n-1})}\\
&\ge\frac{(q^{6n}-q^{2n})(q^{6n}-q^{2n+1})\cdots(q^{6n}-q^{6n-1})}{q^{6n}\cdot q^{6n}\cdots q^{6n}}\\
&=(1-1/q^{4n})(1-1/q^{4n-1})\cdots(1-1/q)\ge1-\sum_{i=1}^{4n}1/q^i\ge1-2/q.
\end{align*}
It follows that with probability $\ge1-4/q$, we can assume in addition that $W_i$ form a direct sum decomposition of $\mathbb F_q^{6n}$.
Therefore, we turn to bound the probability that there exists a non-scalar invertible matrix stabilizing these two direct sum decompositions of $\mathbb F_q^{6n}$. Since $i_k$ are all different, the two direct sum decompositions $U_1\oplus U_2\oplus U_3$ and $W_1\oplus W_2\oplus W_3$ are independent.
So we can assume that $U_i$ is spanned by those standard basis vectors $\vec e_{2n(i-1)+1},\ldots,\vec e_{2ni}$, $i=1,2,3$. The group that stabilizes this direct sum decomposition $U_1\oplus U_2\oplus U_3$ consists of $\begin{bmatrix}D_1&0&0\\0&D_2&0\\0&0&D_3\end{bmatrix}\in\mathrm{GL}(6n,\mathbb F_q)$ where $D_i$ is of size $2n\times2n$.
The question then becomes to bound the probability for a random $W_1\oplus W_2\oplus W_3$ to be stabilized by a non-scalar matrix of the above form. This can be formulated as the following linear algebraic problem. (Recall the correspondence between random $m$-dimensional subspaces and random rank-$m$ matrices as discussed at the beginning of the subsection.) Let $W=\begin{bmatrix}W_{11}&W_{12}&W_{13}\\W_{21}&W_{22}&W_{23}\\W_{31}&W_{32}&W_{33}\end{bmatrix}\in\mathrm{GL}(6n,q)$ be a block matrix where $W_{ij}$ is of size $2n\times2n$. Suppose the columns of $\begin{bmatrix}W_{1i}\\W_{2i}\\W_{3i}\end{bmatrix}$ span $W_i$. Then $D=\operatorname{diag}(D_1,D_2,D_3)$ stabilizes $W_1\oplus W_2\oplus W_3$ if and only if there exists a block diagonal matrix $E=\operatorname{diag}(E_1,E_2,E_3)$, $E_i\in\mathrm{GL}(2n,q)$, such that
\begin{equation}\tag{5}
\begin{bmatrix}D_1&0&0\\0&D_2&0\\0&0&D_3\end{bmatrix}\begin{bmatrix}W_{11}&W_{12}&W_{13}\\W_{21}&W_{22}&W_{23}\\W_{31}&W_{32}&W_{33}\end{bmatrix}=\begin{bmatrix}W_{11}&W_{12}&W_{13}\\W_{21}&W_{22}&W_{23}\\W_{31}&W_{32}&W_{33}\end{bmatrix}\begin{bmatrix}E_1&0&0\\0&E_2&0\\0&0&E_3\end{bmatrix}.
\end{equation}
Note that each direct sum decomposition $W_1\oplus W_2\oplus W_3$, $\dim(W_i)=2n$, has $6\cdot|\mathrm{GL}(2n,q)|^3$ such matrix representations. (The factor 6 takes care of the orders of the three summands.) So the question becomes to bound the probability for a random invertible matrix to have a non-scalar $D$ and $E$ satisfying Equation5.
First, note that Equation5 holds if and only if $D_iW_{i,j}=W_{i,j}E_j$ for $i,j\in[3]$.
\begin{authorclaim}{Claim20}
When $q=\Omega(1)$, we have $\Pr[\forall i,j\in[3],\operatorname{rk}(W_{i,j})=2n]\ge1-20/q$.
\end{authorclaim}
\begin{proof}
Let us work in the setting when $W$ is a random matrix, not necessarily the one above. Then $\Pr[\operatorname{rk}(W)=6n]\ge1-2/q$. For any $i,j\in[3]$, $\Pr[\operatorname{rk}(W_{i,j})<2n]\le2/q$, so $\Pr[\exists i,j\in[3],\operatorname{rk}(W_{i,j})<2n]\le18/q$. It follows that
\begin{align*}
&\Pr[\exists i,j\in[3],\operatorname{rk}(W_{i,j})<2n\mid\operatorname{rk}(W)=6n]\\
&=\frac{\Pr[\exists i,j\in[3],\operatorname{rk}(W_{i,j})<2n\wedge\operatorname{rk}(W)=6n]}{\Pr[\operatorname{rk}(W)=6n]}\\
&\le\frac{18/q}{1-2/q}=\frac{18}{q-2}\le\frac{20}q,
\end{align*}
where the last inequality uses that $q=\Omega(1)$.
\end{proof}
So we assume that $\operatorname{rk}(W_{i,j})=2n$ for all $i,j\in[3]$ in the following, with a loss of probability $\le20/q$.
For $i\in[3]$, by $D_iW_{ii}=W_{ii}E_i$, we have $D_i=W_{ii}E_iW_{ii}^{-1}$. For $i\ne j$, by $(W_{jj}E_jW_{jj}^{-1})W_{ji}=D_jW_{ji}=W_{ji}E_i$, we have $E_j=W_{jj}^{-1}W_{ji}E_iW_{ji}^{-1}W_{jj}$. Again for $i\ne j$, we have $W_{ii}E_iW_{ii}^{-1}W_{ij}=D_iW_{ij}=W_{ij}E_j=W_{ij}W_{jj}^{-1}W_{ji}E_iW_{ji}^{-1}W_{jj}$. It follows that
\[\forall i,j\in[3],i\ne j,\quad E_iW_{ii}^{-1}W_{ij}W_{jj}^{-1}W_{ji}=W_{ii}^{-1}W_{ij}W_{jj}^{-1}W_{ji}E_i.\]
In particular, $E_3$ commutes with $X=W_{33}^{-1}W_{32}W_{22}^{-1}W_{23}$ and $Y=W_{33}^{-1}W_{31}W_{11}^{-1}W_{13}$. Since $W_{ij}$ are independent random invertible matrices, $X$ and $Y$ are independent random invertible matrices. We now resort to the following classical result.
\begin{authorclaim}{Theorem21 ([41], cf. also [23,40])}
Let $X$ and $Y$ be two random matrices in $\mathrm{SL}(n,q)$. Then the probability of $X$ and $Y$ not generating $\mathrm{SL}(n,q)$ is $\le1/q^{\Omega(n)}$.
\end{authorclaim}
Back to our setting, the above theorem implies that the group $G$ generated by random $X$ and $Y$ from $\mathrm{GL}(2n,q)$ contains $\mathrm{SL}(2n,q)$ with probability $\ge1-1/q^{\Omega(n)}$. It follows that $E_3$ belongs to the centralizer of $G$, so $E_3$ must be a scalar matrix. Then note that $D_i$'s and other $E_i$'s are all conjugates of $E_3$. So we have $\forall i\in[3],D_i=E_i=\lambda I_{2n}$ for some $\lambda\ne0\in\mathbb F_q$.
Summarizing the above, we have
\begin{align*}
&\Pr[\mathcal S\text{ non-scalar for }V_1,\ldots,V_6]\\
&\le\Pr[\mathcal S\text{ non-scalar for }V_i\wedge\mathbb F_q^{6n}=U_1\oplus U_2\oplus U_3=W_1\oplus W_2\oplus W_3]+4/q\\
&\le\Pr[\mathcal S\text{ non-scalar for }V_i\mid\mathbb F_q^{6n}=U_1\oplus U_2\oplus U_3=W_1\oplus W_2\oplus W_3]+4/q\\
&\le\Pr[D\text{ non-scalar for }W\wedge\forall i,j\in[3],\operatorname{rk}(W_{ij})=2n]+20/q+4/q\\
&\le\Pr[D\text{ non-scalar for }W\mid\forall i,j\in[3],\operatorname{rk}(W_{ij})=2n]+24/q\\
&\le1/q^{\Omega(n)}+24/q\le1/q^{\Omega(1)},
\end{align*}
when $q=n^{\Omega(1)}$. This concludes the proof of Proposition15. \hfill$\square$

\section*{6.2 Constructive Baer Correspondence and Theorems A and B}
Let us review Baer's Correspondence [7], which connects alternating bilinear maps with $p$-groups of class 2 and exponent $p$. Let $P$ be a $p$-group of class 2 and exponent $p$, $p>2$. Suppose the commutator subgroup $[P,P]\cong\mathbb Z_p^m$ and $P/[P,P]\cong\mathbb Z_p^n$. Then the commutator map $[,]:P/[P,P]\times P/[P,P]\to[P,P]$ is an alternating bilinear map. Conversely, let $\phi:\mathbb Z_p^n\times\mathbb Z_p^n\to\mathbb Z_p^m$ be an alternating bilinear map. Then a $p$-group of class 2 and exponent $p$, denoted as $P_\phi$ can be defined as follows. The group elements are from $\mathbb Z_p^n\times\mathbb Z_p^m$, and the group product $\cdot$ is defined as
\[(u,v)\cdot(u',v')=(u+u',v+v'+\tfrac12\phi(u,u')).\]
We say that $(A,B)\in\mathrm{GL}(n,p)\times\mathrm{GL}(m,p)$ is a pseudo-autometry of $\phi$, if $\phi=B\circ\phi\circ A$. Wilson [71] elucidated the structure of $\operatorname{Aut}(P_\phi)$ in terms of the pseudo-autometry group of $\phi$, that we denote $\Psi\operatorname{Aut}(\phi)$. Here we recall the consequence of Wilson's result that we need for counting group isomorphisms.
\begin{authorclaim}{Proposition23 (Wilson [71, Prop.3.8], see [15, Prop.2.4] for notation closer to ours)}
For $\phi:\mathbb Z_p^n\times\mathbb Z_p^n\to\mathbb Z_p^m$ an alternating bilinear map,
\[|\operatorname{Aut}(P_\phi)|=|\Psi\operatorname{Aut}(\phi)|p^{nm},\]
where $\Psi\operatorname{Aut}(\phi)$ denotes the pseudo-autometry group of $\phi$.
\end{authorclaim}
We then state a lemma which can be viewed as a constructive version of Baer's Correspondence, communicated to us by James B. Wilson.
\begin{authorclaim}{Lemma24 (Constructive version of Baer's Correspondence for matrix groups)}
Let $p$ be an odd prime. Over the finite field $\mathbb F=\mathbb F_{p^e}$, Alternating Matrix Space Isometry is equivalent to Group Isomorphism for matrix groups over $\mathbb F$ that are $p$-groups of class 2 and exponent $p$. More precisely, there are functions computable in time $\operatorname{poly}(n,m,\log|\mathbb F|)$:
\begin{itemize}
\item $G:\Lambda(n,\mathbb F)^m\to\mathrm M(n+m+1,\mathbb F)^{n+m}$ and
\item $\operatorname{Alt}:\mathrm M(n,\mathbb F)^m\to\Lambda(m,\mathbb F)^{O(m^2)}$
\end{itemize}
such that: (1) for an alternating bilinear map $\mathbf A$, the group generated by $G(\mathbf A)$ is the Baer group corresponding to $\mathbf A$, (2) $G$ and Alt are mutually inverse, in the sense that the group generated by $G(\operatorname{Alt}(M_1,\ldots,M_m))$ is isomorphic to the group generated by $M_1,\ldots,M_m$, and conversely $\operatorname{Alt}(G(\mathbf A))$ is pseudo-isometric to $\mathbf A$.
\end{authorclaim}
\begin{proof}
First, let $G$ be a $p$-group of class 2 and exponent $p$ given by $m$ generating matrices of size $n\times n$ over $\mathbb F$. Then from the generating matrices of $G$, we first compute a generating set of $[G,G]$, by just computing all the commutators of the given generators. We can then remove those redundant elements from this generating set in time $\operatorname{poly}(\log|[G,G]|,\log|\mathbb F|)$, using Luks' result on computing with solvable matrix groups [51]. We then compute a set of representatives of a non-redundant generating set of $G/[G,G]$, again using Luks's aforementioned result. From these data we can compute an alternating bilinear map representing the commutator map of $G$ in time $\operatorname{poly}(n,m,\log|\mathbb F|)$.
Conversely, let an alternating bilinear map be given by $\mathbf A=(A_1,\ldots,A_m)\in\Lambda(n,\mathbb F)^m$. From $\mathbf A$, for $i\in[n]$, construct $B_i=[A_1e_i,\ldots,A_me_i]\in\mathrm M(n\times m,\mathbb F)$, where $e_i$ is the $i$th standard basis vector of $\mathbb F^n$. That is, the $j$th column of $B_i$ is the $i$th column of $A_j$. Then for $i\in[n]$, construct
\[\widetilde B_i=\begin{bmatrix}1&e_i^t&0\\0&I_n&B_i\\0&0&I_m\end{bmatrix}\in\mathrm{GL}(1+n+m,\mathbb F),\]
where $e_i\in\mathbb F^n$, and for $j\in[m]$, construct
\[\widetilde C_j=\begin{bmatrix}1&0&e_j^t\\0&I_n&0\\0&0&I_m\end{bmatrix}\in\mathrm{GL}(1+n+m,\mathbb F),\]
where $e_j\in\mathbb F^m$. Let $G(\mathbf A)$ be the matrix group generated by $\widetilde B_i$ and $\widetilde C_j$. Then it can be verified easily that, $G(\mathbf A)$ is isomorphic to the Baer group corresponding to the alternating bilinear map defined by $\mathbf A$. In particular, $[G,G]\cong\mathbb F^m\cong\mathbb Z_p^{em}$ (isomorphism of abelian groups), and $G/[G,G]\cong\mathbb F^n\cong\mathbb Z_p^{en}$. This construction can be done in time $\operatorname{poly}(n,m,\log|\mathbb F|)$.
\end{proof}

\begin{proof}[Proof of TheoremB]
For the counting-to-decision reduction, we basically follow the above routine, but with a twist, because of the minor distinction between alternating matrix space isometry, and alternating bilinear map pseudo-isometry. Let us briefly explain this issue. Suppose from an alternating bilinear map $\phi:\mathbb Z_p^n\times\mathbb Z_p^n\to\mathbb Z_p^m$ we constructed the $p$-group $P_\phi$ of class 2 and exponent $p$; by Proposition23 $|\operatorname{Aut}(P_\phi)|=p^{nm}|\Psi\operatorname{Aut}(\phi)|$, so by multiplying the result by $p^{nm}$, it is necessary and sufficient to count the psuedo-autometries of $\phi$.
Towards that end, let $(C_1,\ldots,C_m)\in\Lambda(n,p)$ be a matrix representation of $\phi$. If $C_i$'s are linearly independent, then for a pseudo-autometry $(A,B)\in\mathrm{GL}(n,p)\times\mathrm{GL}(m,p)$, given $A$ there exists a unique $B$ that makes $(A,B)$ a pseudo-autometry. If $C_i$'s are not linearly independent, say the linear span of $C_i$'s is of dimension $m'$, then the number of $B$ such that $(A,B)$ is a pseudo-autometry (assuming there are any) is $|\mathrm M((m-m')\times m',p)||\mathrm{GL}(m-m',p)|=p^{m'(m-m')}|\mathrm{GL}(m-m',p)|$. To see this, suppose that we have taken linear combinations of the $C_i$ so that $C_1,\ldots,C_{m'}$ are linearly independent and $C_{m'+1},C_{m'+2},\ldots,C_m$ are zero. Then without changing the $C_i$, we may take any invertible linear combination among $C_{m'+1},\ldots,C_m$ (a copy of $\mathrm{GL}(m-m',p)$), and we may add any linear combination of the last $m-m'$ matrices to the first $m'$ matrices (a copy of $\mathrm M((m-m')\times m',p)$). The counting to decision reduction for Alternating Matrix Space Isometry computes the number of $A\in\mathrm{GL}(n,p)$ so that there exists some $B\in\mathrm{GL}(m,p)$ such that $(A,B)$ is a pseudo-autometry. So it needs to be multiplied by a factor of $p^{m'(m-m')}|\mathrm{GL}(m-m',p)|$.
\end{proof}

\begin{thebibliography}{74}

\bibitem[1]{ref1}
Eric Allender and Bireswar Das. Zero knowledge and circuit minimization. Inf. Comput., 256:2–8, 2017. doi:10.1016/j.ic.2017.04.004.

\bibitem[2]{ref2}
Vikraman Arvind and Jacobo Torán. Isomorphism testing: Perspective and open problems. Bulletin of the EATCS, 86:66–84, 2005.

\bibitem[3]{ref3}
László Babai. Graph isomorphism in quasipolynomial time [extended abstract]. In Proceedings of the 48th Annual ACM SIGACT Symposium on Theory of Computing, STOC 2016, pages 684–697, 2016. arXiv:1512.03547 [cs.DS] version 2. doi:10.1145/2897518.2897542.

\bibitem[4]{ref4}
László Babai, Paolo Codenotti, Joshua A. Grochow, and Youming Qiao. Code equivalence and group isomorphism. In Proceedings of the Twenty-Second Annual ACM–SIAM Symposium on Discrete Algorithms (SODA11), pages 1395–1408, Philadelphia, PA, 2011. SIAM. doi: 10.1137/1.9781611973082.107.

\bibitem[5]{ref5}
László Babai, Paolo Codenotti, and Youming Qiao. Polynomial-time isomorphism test for groups with no abelian normal subgroups - (extended abstract). In Automata, Languages, and Programming - 39th International Colloquium, ICALP 2012, Proceedings, Part I, pages 51–62, 2012. doi:10.1007/978-3-642-31594-7\_5.

\bibitem[6]{ref6}
László Babai and Youming Qiao. Polynomial-time isomorphism test for groups with Abelian Sylow towers. In 29th STACS, pages 453–464. Springer LNCS 6651, 2012. doi:10.4230/ LIPIcs.STACS.2012.453.

\bibitem[7]{ref7}
Reinhold Baer. Groups with abelian central quotient group. Trans. AMS, 44(3):357–386, 1938. doi:10.1090/S0002-9947-1938-1501972-1.

\bibitem[8]{ref8}
Mihir Bellare and Shafi Goldwasser. The complexity of decision versus search. SIAM J. Comput., 23(1):97–119, 1994. doi:10.1137/S0097539792228289.

\bibitem[9]{ref9}
Hans Ulrich Besche and Bettina Eick. Construction of finite groups. J. Symb. Comput., 27(4):387–404, 1999. doi:10.1006/jsco.1998.0258.

\bibitem[10]{ref10}
Hans Ulrich Besche, Bettina Eick, and E.A. O’Brien. A millennium project: Constructing small groups. Intern. J. Alg. and Comput, 12:623–644, 2002. doi:10.1142/S0218196702001115.

\bibitem[11]{ref11}
Anton Betten, Michael Braun, Harald Fripertinger, Adalbert Kerber, Axel Kohnert, and Alfred Wassermann. Error-correcting linear codes: Classification by isometry and applications, volume 18. Springer Science and Business Media, 2006.

\bibitem[12]{ref12}
Peter A. Brooksbank, Joshua A. Grochow, Yinan Li, Youming Qiao, and James B. Wilson. Incorporating Weisfeiler–Leman into algorithms for group isomorphism. arXiv:1905.02518 [cs.CC], 2019.

\bibitem[13]{ref13}
Peter A. Brooksbank, Yinan Li, Youming Qiao, and James B. Wilson. Improved algorithms for alternating matrix space isometry: From theory to practice. In Fabrizio Grandoni, Grzegorz Herman, and Peter Sanders, editors, 28th Annual European Symposium on Algorithms, ESA 2020, September 7-9, 2020, Pisa, Italy (Virtual Conference), volume 173 of LIPIcs, pages 26:1–26:15. Schloss Dagstuhl - Leibniz-Zentrum für Informatik, 2020. doi:10.4230/LIPIcs. ESA.2020.26.

\bibitem[14]{ref14}
Peter A. Brooksbank and Eugene M. Luks. Testing isomorphism of modules. J. Algebra, 320(11):4020–4029, 2008. doi:10.1016/j.jalgebra.2008.07.014.

\bibitem[15]{ref15}
Peter A. Brooksbank, Joshua Maglione, and James B. Wilson. A fast isomorphism test for groups whose Lie algebra has genus 2. J. Algebra, 473:545–590, 2017. doi:10.1016/j. jalgebra.2016.12.007.

\bibitem[16]{ref16}
Peter A. Brooksbank and E. A. O’Brien. Constructing the group preserving a system of forms. Internat. J. Algebra Comput., 18(2):227–241, 2008. doi:10.1142/S021819670800441X.

\bibitem[17]{ref17}
John J. Cannon and Derek F. Holt. Automorphism group computation and isomorphism testing in finite groups. J. Symbolic Comput., 35(3):241–267, 2003. doi:10.1016/S0747-7171(02) 00133-5.

\bibitem[18]{ref18}
Xi Chen, Xiaotie Deng, and Shang-Hua Teng. Settling the complexity of computing two-player Nash equilibria. J. ACM, 56(3):Art. 14, 57, 2009. doi:10.1145/1516512.1516516.

\bibitem[19]{ref19}
Alexander Chistov, Gábor Ivanyos, and Marek Karpinski. Polynomial time algorithms for modules over finite dimensional algebras. In Proceedings of the 1997 International Symposium on Symbolic and Algebraic Computation, ISSAC ’97, pages 68–74. ACM, 1997. doi:10.1145/ 258726.258751.

\bibitem[20]{ref20}
Serena Cicalò, Willem A. de Graaf, and Michael Vaughan-Lee. An effective version of the Lazard correspondence. J. Algebra, 352(1):430–450, 2012. doi:10.1016/j.jalgebra.2011.11.031.

\bibitem[21]{ref21}
W.A. de Graaf. Lie Algebras: Theory and Algorithms, volume 56 of North-Holland Mathemat- ical Library. Elsevier Science, 2000.

\bibitem[22]{ref22}
Holger Dell, Thore Husfeldt, Dániel Marx, Nina Taslaman, and Martin Wahlén. Exponential time complexity of the permanent and the Tutte polynomial. ACM Trans. Algorithms, 10(4):Art. 21, 32, 2014. doi:10.1145/2635812.

\bibitem[23]{ref23}
Sean Eberhard and Stefan-C. Virchow. Random generation of the special linear group. Transactions of the American Mathematical Society, page 1, 2020. doi:10.1090/tran/8009.

\bibitem[24]{ref24}
Wayne Eberly and Mark Giesbrecht. Efficient decomposition of associative algebras over finite fields. Journal of Symbolic Computation, 29(3):441–458, 2000. doi:10.1006/jsco.1999.0308.

\bibitem[25]{ref25}
Bettina Eick, C. R. Leedham-Green, and E. A. O’Brien. Constructing automorphism groups of p-groups. Comm. Algebra, 30(5):2271–2295, 2002. doi:10.1081/AGB-120003468.

\bibitem[26]{ref26}
V. Felsch and J. Neubüser. On a programme for the determination of the automorphism group of a finite group. In Pergamon J. Leech, editor, Computational Problems in Abstract Algebra (Proceedings of a Conference on Computational Problems in Algebra, Oxford, 1967), pages 59–60, Oxford, 1970.

\bibitem[27]{ref27}
Katalin Friedl and Lajos Rónyai. Polynomial time solutions of some problems in computational algebra. In Robert Sedgewick, editor, Proceedings of the 17th Annual ACM Symposium on Theory of Computing, May 6-8, 1985, Providence, Rhode Island, USA, pages 153–162. ACM, 1985. doi:10.1145/22145.22162.

\bibitem[28]{ref28}
Vyacheslav Futorny, Joshua A. Grochow, and Vladimir V. Sergeichuk. Wildness for tensors. Lin. Alg. Appl., 566:212–244, 2019. doi:10.1016/j.laa.2018.12.022.

\bibitem[29]{ref29}
Joshua A. Grochow. Answer to “what is the hardest instance for the group isomorph- ism problem?”. Theoretical Computer Science Stack Exchange. URL: https://cstheory. stackexchange.com/a/42551/129.

\bibitem[30]{ref30}
Joshua A. Grochow and Youming Qiao. Polynomial-time isomorphism test of groups that are tame extensions - (extended abstract). In Algorithms and Computation - 26th International Symposium, ISAAC 2015, Nagoya, Japan, December 9-11, 2015, Proceedings, pages 578–589, 2015. doi:10.1007/978-3-662-48971-0\_49.

\bibitem[31]{ref31}
Joshua A. Grochow and Youming Qiao. Algorithms for group isomorphism via group extensions and cohomology. SIAM J. Comput., 46(4):1153–1216, 2017. Preliminary version in IEEE Conference on Computational Complexity (CCC) 2014 (DOI:10.1109/CCC.2014.19). Also available as arXiv:1309.1776 [cs.DS] and ECCC Technical Report TR13-123. doi:10.1137/ 15M1009767.

\bibitem[32]{ref32}
Joshua A. Grochow and Youming Qiao. On the complexity of isomorphism problems for tensors, groups, and polynomials I: Tensor Isomorphism-completeness. In ITCS, page to appear, 2021. arXiv:1907.00309.

\bibitem[33]{ref33}
Martin Grohe and Pascal Schweitzer. The graph isomorphism problem. Commun. ACM, 63(11):128–134, 2020. doi:10.1145/3372123.

\bibitem[34]{ref34}
Xiaoyu He and Youming Qiao. On the Baer–Lovász–Tutte construction of groups from graphs: isomorphism types and homomorphism notions. arXiv:2003.07200 [math.CO], 2020.

\bibitem[35]{ref35}
Russell Impagliazzo and Ramamohan Paturi. On the complexity of k-SAT. J. Comput. System Sci., 62(2):367–375, 2001. Special issue on the Fourteenth Annual IEEE Conference on Computational Complexity (Atlanta, GA, 1999). doi:10.1006/jcss.2000.1727.

\bibitem[36]{ref36}
Gábor Ivanyos. Fast randomized algorithms for the structure of matrix algebras over finite fields. In Proceedings of the 2000 international symposium on Symbolic and algebraic computation, pages 175–183. ACM, 2000. doi:10.1145/345542.345620.

\bibitem[37]{ref37}
Gábor Ivanyos, Marek Karpinski, and Nitin Saxena. Deterministic polynomial time algorithms for matrix completion problems. SIAM J. Comput., 39(8):3736–3751, 2010. doi:10.1137/ 090781231.

\bibitem[38]{ref38}
Gábor Ivanyos and Lajos Rónyai. Computations in associative and Lie algebras. In Some tapas of computer algebra, pages 91–120. Springer, 1999. doi:10.1007/978-3-662-03891-8\_5.

\bibitem[39]{ref39}
Zhengfeng Ji, Youming Qiao, Fang Song, and Aaram Yun. General linear group action on tensors: A candidate for post-quantum cryptography. In Dennis Hofheinz and Alon Rosen, editors, Theory of Cryptography - 17th International Conference, TCC 2019, Nuremberg, Germany, December 1-5, 2019, Proceedings, Part I, volume 11891 of Lecture Notes in Computer Science, pages 251–281. Springer, 2019. Preprint arXiv:1906.04330 [cs.CR]. doi:10.1007/ 978-3-030-36030-6\_11.

\bibitem[40]{ref40}
William M. Kantor. Some topics in asymptotic group theory. Groups, Combinatorics and Geometry (Durham, pages 403–421, 1990.

\bibitem[41]{ref41}
William M Kantor and Alexander Lubotzky. The probability of generating a finite classical group. Geometriae Dedicata, 36(1):67–87, 1990.

\bibitem[42]{ref42}
Neeraj Kayal and Timur Nezhmetdinov. Factoring groups efficiently. In Susanne Albers, Alberto Marchetti-Spaccamela, Yossi Matias, Sotiris E. Nikoletseas, and Wolfgang Thomas, editors, Automata, Languages and Programming, 36th International Colloquium, ICALP 2009, Rhodes, Greece, July 5-12, 2009, Proceedings, Part I, volume 5555 of Lecture Notes in Computer Science, pages 585–596. Springer, 2009. Preprint ECCC Tech. Report TR08-074. doi:10.1007/978-3-642-02927-1\_49.

\bibitem[43]{ref43}
E. I. Khukhro. p-automorphisms of finite p-groups, volume 246 of London Mathematical Society Lecture Note Series. Cambridge University Press, Cambridge, 1998. doi:10.1017/ CBO9780511526008.

\bibitem[44]{ref44}
Johannes Köbler, Uwe Schöning, and Jacobo Torán. The graph isomorphism problem: its structural complexity. Birkhauser Verlag, Basel, Switzerland, Switzerland, 1993. doi:10.1007/ 978-1-4612-0333-9.

\bibitem[45]{ref45}
Tamara G Kolda and Brett W Bader. Tensor decompositions and applications. SIAM review, 51(3):455–500, 2009. doi:10.1137/07070111X.

\bibitem[46]{ref46}
Michel Lazard. Sur les groupes nilpotents et les anneaux de Lie. Ann. Sci. Ecole Norm. Sup. (3), 71:101–190, 1954. doi:10.24033/asens.1021.

\bibitem[47]{ref47}
François Le Gall. Efficient isomorphism testing for a class of group extensions. In Proc. 26th STACS, pages 625–636, 2009. doi:10.4230/LIPIcs.STACS.2009.1830.

\bibitem[48]{ref48}
Mark L. Lewis and James B. Wilson. Isomorphism in expanding families of indistinguishable groups. Groups Complex. Cryptol., 4(1):73–110, 2012. doi:10.1515/gcc-2012-0008.

\bibitem[49]{ref49}
Yinan Li and Youming Qiao. Linear algebraic analogues of the graph isomorphism problem and the Erdős–Rényi model. In Chris Umans, editor, 58th IEEE Annual Symposium on Foundations of Computer Science, FOCS 2017, pages 463–474. IEEE Computer Society, 2017. doi:10.1109/FOCS.2017.49.

\bibitem[50]{ref50}
Richard J. Lipton, Lawrence Snyder, and Yechezkel Zalcstein. The complexity of word and isomorphism problems for finite groups. Yale University Department of Computer Science Research Report \# 91, 1977. URL: https://apps.dtic.mil/dtic/tr/fulltext/u2/a053246. pdf.

\bibitem[51]{ref51}
Eugene M. Luks. Computing in solvable matrix groups. In FOCS 1992, 33rd Annual Symposium on Foundations of Computer Science, pages 111–120. IEEE Computer Society, 1992. doi:10.1109/SFCS.1992.267813.

\bibitem[52]{ref52}
Eugene M. Luks. Permutation groups and polynomial-time computation. In Groups and computation (New Brunswick, NJ, 1991), volume 11 of DIMACS Ser. Discrete Math. Theoret. Comput. Sci., pages 139–175. Amer. Math. Soc., Providence, RI, 1993.

\bibitem[53]{ref53}
Eugene M. Luks. Hypergraph isomorphism and structural equivalence of boolean functions. In Proceedings of the Thirty-First Annual ACM Symposium on Theory of Computing, May 1-4, 1999, Atlanta, Georgia, USA, pages 652–658, 1999. doi:10.1145/301250.301427.

\bibitem[54]{ref54}
Rudolf Mathon. A note on the graph isomorphism counting problem. Information Processing Letters, 8(3):131–136, 1979.

\bibitem[55]{ref55}
Brendan D. McKay. Practical graph isomorphism. Congr. Numer., pages 45–87, 1980.

\bibitem[56]{ref56}
Brendan D. McKay and Adolfo Piperno. Practical graph isomorphism, II. Journal of Symbolic Computation, 60(0):94–112, 2014. doi:10.1016/j.jsc.2013.09.003.

\bibitem[57]{ref57}
Alan H. Mekler. Stability of nilpotent groups of class 2 and prime exponent. The Journal of Symbolic Logic, 46(4):781–788, 1981.

\bibitem[58]{ref58}
Gary L. Miller. On the $n^{\log n}$ isomorphism technique (a preliminary report). In STOC, pages 51–58. ACM, 1978. doi:10.1145/800133.804331.

\bibitem[59]{ref59}
Takunari Miyazaki. Luks’s reduction of graph isomorphism to code equivalence. Comment to E. W. Clark, https://groups.google.com/forum/\#!msg/sci.math.research/puZxGj9HXKI/ CeyH2yyyNFUJ, 1996.

\bibitem[60]{ref60}
Vipul Naik. Lazard correspondence up to isoclinism. PhD thesis, The University of Chicago, 2013. URL: https://vipulnaik.com/thesis/.

\bibitem[61]{ref61}
Jacques Patarin. Hidden fields equations (HFE) and isomorphisms of polynomials (IP): two new families of asymmetric algorithms. In Advances in Cryptology - EUROCRYPT ’96, Inter- national Conference on the Theory and Application of Cryptographic Techniques, Saragossa, Spain, May 12-16, 1996, Proceeding, pages 33–48, 1996. doi:10.1007/3-540-68339-9\_4.

\bibitem[62]{ref62}
Erez Petrank and Ron M. Roth. Is code equivalence easy to decide? IEEE Trans. Inf. Theory, 43(5):1602–1604, 1997. doi:10.1109/18.623157.

\bibitem[63]{ref63}
Youming Qiao, Jayalal M. N. Sarma, and Bangsheng Tang. On isomorphism testing of groups with normal Hall subgroups. In Proc. 28th STACS, pages 567–578, 2011. doi: 10.4230/LIPIcs.STACS.2011.567.

\bibitem[64]{ref64}
Lajos Rónyai. Computing the structure of finite algebras. J. Symb. Comput., 9(3):355–373, 1990. doi:10.1016/S0747-7171(08)80017-X.

\bibitem[65]{ref65}
David J. Rosenbaum. Bidirectional collision detection and faster deterministic isomorphism testing. arXiv preprint arXiv:1304.3935 [cs.DS], 2013.

\bibitem[66]{ref66}
David J. Rosenbaum. Breaking the $n^{\log n}$ barrier for solvable-group isomorphism. In Proceedings of the Twenty-Fourth Annual ACM-SIAM Symposium on Discrete Algorithms, pages 1054–1073. SIAM, 2013. Preprint arXiv:1205.0642 [cs.DS]. doi:10.1137/1.9781611973105.76.

\bibitem[67]{ref67}
Nicolas Sendrier and Dimitris E. Simos. The hardness of code equivalence over $\mathbb F_q$ and its application to code-based cryptography. In International Workshop on Post-Quantum Cryptography, pages 203–216. Springer, 2013.

\bibitem[68]{ref68}
Seinosuke Toda. PP is as hard as the polynomial-time hierarchy. SIAM J. Comput., 20(5):865– 877, 1991. doi:10.1137/0220053.

\bibitem[69]{ref69}
Leslie G. Valiant. Relative complexity of checking and evaluating. Information Processing Lett., 5(1):20–23, 1976/77. doi:10.1016/0020-0190(76)90097-1.

\bibitem[70]{ref70}
James Wilson. 2014 conference on Groups, Computation, and Geometry at Colorado State University, co-organized by P. Brooksbank, A. Hulpke, T. Penttila, J. Wilson, and W. Kantor. Personal communication, 2014.

\bibitem[71]{ref71}
James B. Wilson. Decomposing p-groups via Jordan algebras. J. Algebra, 322(8):2642–2679, 2009. doi:10.1016/j.jalgebra.2009.07.029.

\bibitem[72]{ref72}
James B. Wilson. Finding direct product decompositions in polynomial time. arXiv:1005.0548 [math.GR], 2010.

\bibitem[73]{ref73}
James B. Wilson. Existence, algorithms, and asymptotics of direct product decompositions, I. Groups Complex. Cryptol., 4(1):33–72, 2012. doi:10.1515/gcc-2012-0007.

\bibitem[74]{ref74}
V. N. Zemlyachenko, N. M. Korneenko, and R. I. Tyshkevich. Graph isomorphism problem. J. Soviet Math., 29(4):1426–1481, May 1985. doi:10.1007/BF02104746.
\end{thebibliography}
\end{document}
